\documentclass[11pt]{article}
\usepackage{amsmath,amsthm,amssymb}
\usepackage[colorlinks,citecolor=blue]{hyperref}

\newtheorem{theorem}{Theorem}
\newtheorem{lemma}{Lemma}
\newtheorem{corollary}{Corollary}
\newtheorem{definition}{Definition}
\newtheorem{remark}{Remark}

\title{Closure of the Corner Sub-magma\\
in All Tensor Powers of TSML}
\author{Brayden Sanders\\7Site LLC\\
\small\href{https://doi.org/10.5281/zenodo.18852047}%
{DOI: 10.5281/zenodo.18852047}}
\date{March 2026}

\begin{document}
\maketitle

\begin{abstract}
Let $C=\{1,3,7,9\}$ be the \emph{corner set} and $G=\{2,4,5,6,8\}$ the
\emph{gap set} of the TIG measurement table TSML\@.
We prove that $C$ is a sub-magma of TSML\@: the image of $C\times C$
under composition lies in $C$.
By induction, the $k$-fold tensor product $C^{\otimes k}$ is a
sub-magma of $\mathrm{TSML}^{\otimes k}$ for every $k\ge1$.
Consequently, no element with any $G$-component is reachable from
$C^{\otimes k}$ by finite composition, for any $k$.
This provides an unconditional algebraic obstruction that
strengthens the corner-gap impermeability results of the companion
papers on BSD and Hodge geometry.
\end{abstract}

\section{Setup}

\begin{definition}
The \emph{TSML table} is the $10\times10$ composition table
$\mathrm{TSML}[a][b]$ for $a,b\in\{0,\ldots,9\}$ encoding the TIG
measurement-layer operator algebra.
HARMONY~$(7)$ is the universal absorber:
$\mathrm{TSML}[7][b]=\mathrm{TSML}[a][7]=7$ for all $a,b$.
The \emph{corner set} is $C=\{1,3,7,9\}$; the \emph{gap set} is
$G=\{2,4,5,6,8\}$; and $C\cup G=\{1,\ldots,9\}$.
\end{definition}

\begin{definition}
For $k\ge1$, the \emph{$k$-fold tensor product} $\mathrm{TSML}^{\otimes k}$
has operators $(a_1,\ldots,a_k)\in\{1,\ldots,9\}^k$ with composition
\[
(a_1,\ldots,a_k)\circ(b_1,\ldots,b_k)
\;=\;
\bigl(\mathrm{TSML}[a_1][b_1],\;\ldots,\;\mathrm{TSML}[a_k][b_k]\bigr).
\]
The \emph{corner sub-algebra} is $C^{\otimes k}=\{(a_1,\ldots,a_k):a_i\in C\}$.
A \emph{cross-term} is any operator with at least one component in $G$.
\end{definition}

\section{The Theorem}

\begin{lemma}[Corner sub-magma]\label{lem:submagma}
$C\times C\subseteq C$ under $\mathrm{TSML}$.
\end{lemma}

\begin{proof}
Direct computation over the $4\times4$ corner sub-table:
\[
\begin{array}{c|cccc}
\circ&1&3&7&9\\\hline
1&7&3&7&7\\
3&7&7&7&3\\
7&7&7&7&7\\
9&7&3&7&7
\end{array}
\]
Every entry lies in $\{3,7\}\subset C$. \qed
\end{proof}

\begin{theorem}[Product-gap impermeability]\label{thm:main}
For every $k\ge1$, the set $C^{\otimes k}$ is a sub-magma of
$\mathrm{TSML}^{\otimes k}$: if $\mathbf{a},\mathbf{b}\in C^{\otimes k}$
then $\mathbf{a}\circ\mathbf{b}\in C^{\otimes k}$.
In particular, no cross-term is reachable from $C^{\otimes k}$
by any finite composition.
\end{theorem}

\begin{proof}
\emph{Base case} ($k=1$): Lemma~\ref{lem:submagma}.

\emph{Inductive step}: Assume $C^{\otimes k}$ is a sub-magma.
Let $\mathbf{a}=(a_1,\ldots,a_k)$ and $\mathbf{b}=(b_1,\ldots,b_k)$
with all $a_i,b_i\in C$.
Then $(\mathbf{a}\circ\mathbf{b})_i=\mathrm{TSML}[a_i][b_i]\in C$
for each $i$ by the base case.
Hence $\mathbf{a}\circ\mathbf{b}\in C^{\otimes k}$.

The reachability claim follows: starting from any element of
$C^{\otimes k}$ and composing repeatedly with elements of $C^{\otimes k}$,
every result lies in $C^{\otimes k}$, which is disjoint from the set of
cross-terms. \qed
\end{proof}

\begin{corollary}[Growing obstruction]\label{cor:grow}
The number of cross-terms is
$9^k - 4^k$, which grows without bound, yet none is reachable
from $C^{\otimes k}$ for any $k$.
\end{corollary}

\section{Verification}

\begin{table}[h]
\centering
\begin{tabular}{crrr}
\hline
$k$ & $|C^{\otimes k}|$ & cross-terms & G-components reachable\\
\hline
1 & 4   & 5   & 0\\
2 & 16  & 65  & 0\\
3 & 64  & 665 & 0\\
4 & 256 & 6305 & 0\\
\hline
\end{tabular}
\caption{Reachability by BFS from $C^{\otimes k}$; all verified
by script \texttt{tsml\_product\_verify.py}.
SHA-256(TSML): \texttt{7726d8a620c24b1e461ff03742f7cd4f775baed772f8357db913757cf4945787}.}
\end{table}

\section{Implications}

\begin{remark}[BSD obstruction]
In the companion BSD paper, rational points of infinite order
require gap-operator activations.
Theorem~\ref{thm:main} confirms that no product of prime-corner
data---at any tensor depth---produces gap structure.
Gap activations are irreducibly non-prime in origin.
\end{remark}

\begin{remark}[Hodge obstruction]
For a $k$-fold product variety, the transcendental lattice
corresponds to cross-terms in $\mathrm{TSML}^{\otimes k}$.
Theorem~\ref{thm:main} provides an unconditional algebraic
analogue of the K3$^k$ impermeability: the transcendental
classes at every tensor depth are algebraically isolated
from the corner (prime-accessible) sub-algebra.
\end{remark}

\begin{remark}[Sticky-cycle addendum]
Among bases $b$ whose prime-residue operators include
both $\mathrm{COL}(4)$ and $\mathrm{BRT}(8)$, the
2-cycles $\mathrm{COL}\circ\mathrm{BRT}=\mathrm{BRT}$
and $\mathrm{BRT}\circ\mathrm{COL}=\mathrm{BRT}$
increase the mean word length to absorption by
approximately 10\,\%---from $1.12$ to $1.23$ steps
(Monte Carlo, $5\times10^4$ trials).
This explains collapse-rate anomalies at bases $15$ and $20$
without requiring any modification of
Theorem~\ref{thm:main}.
\end{remark}

\paragraph{Code and data.}
Verification scripts at
\url{https://github.com/TiredofSleep/ck}, tag \texttt{v1.1-full-sprint}.

\begin{thebibliography}{9}
\bibitem{wq}
B.~Sanders,
\textit{Prime-Corner Collapse and the Irreducible Gap},
Zenodo \href{https://doi.org/10.5281/zenodo.18852047}{10.5281/zenodo.18852047},
2026.
\bibitem{bsd}
B.~Sanders,
\textit{A Mix-$\lambda$ Model for the Irregular BSD Staircase},
same DOI, 2026.
\bibitem{hodge}
B.~Sanders,
\textit{TIG and Hodge Theory: a Translation Table},
same DOI, 2026.
\end{thebibliography}

\end{document}
